\documentclass[10pt,a4paper]{amsart}
\usepackage{amsmath,amssymb,amsthm,graphicx}
\usepackage{verbatim}
\newenvironment{reference}{\comment}{\endcomment}
%\newenvironment{reference}{}{}
\newenvironment{slogan}{\comment}{\endcomment}
\newenvironment{history}{\comment}{\endcomment}
% For commutative diagrams we use Xy-pic
\usepackage[all]{xy}

% We use 2cell for 2-commutative diagrams.
\xyoption{2cell}
% Theorem environments.
%
\theoremstyle{plain}
\newtheorem{theorem}[subsection]{Theorem}
\newtheorem{proposition}[subsection]{Proposition}
\newtheorem{lemma}[subsection]{Lemma}

\theoremstyle{definition}
\newtheorem{definition}[subsection]{Definition}
\newtheorem{example}[subsection]{Example}
\newtheorem{exercise}[subsection]{Exercise}
\newtheorem{situation}[subsection]{Situation}

\theoremstyle{remark}
\newtheorem{remark}[subsection]{Remark}
\newtheorem{remarks}[subsection]{Remarks}

\numberwithin{equation}{subsection}

% Macros
%
\def\lim{\mathop{\mathrm{lim}}\nolimits}
\def\colim{\mathop{\mathrm{colim}}\nolimits}
\def\Spec{\mathop{\mathrm{Spec}}}
\def\Hom{\mathop{\mathrm{Hom}}\nolimits}
\def\Ext{\mathop{\mathrm{Ext}}\nolimits}
\def\SheafHom{\mathop{\mathcal{H}\!\mathit{om}}\nolimits}
\def\SheafExt{\mathop{\mathcal{E}\!\mathit{xt}}\nolimits}
\def\Sch{\mathit{Sch}}
\def\Mor{\mathop{\mathrm{Mor}}\nolimits}
\def\Ob{\mathop{\mathrm{Ob}}\nolimits}
\def\Sh{\mathop{\mathit{Sh}}\nolimits}
\def\NL{\mathop{N\!L}\nolimits}
\def\CH{\mathop{\mathrm{CH}}\nolimits}
\def\proetale{{pro\text{-}\acute{e}tale}}
\def\etale{{\acute{e}tale}}
\def\QCoh{\mathit{QCoh}}
\def\Ker{\mathop{\mathrm{Ker}}}
\def\Im{\mathop{\mathrm{Im}}}
\def\Coker{\mathop{\mathrm{Coker}}}
\def\Coim{\mathop{\mathrm{Coim}}}

% Boxtimes
%
\DeclareMathSymbol{\boxtimes}{\mathbin}{AMSa}{"02}

%
% Macros for moduli stacks/spaces
%
\def\QCohstack{\mathcal{QC}\!\mathit{oh}}
\def\Cohstack{\mathcal{C}\!\mathit{oh}}
\def\Spacesstack{\mathcal{S}\!\mathit{paces}}
\def\Quotfunctor{\mathrm{Quot}}
\def\Hilbfunctor{\mathrm{Hilb}}
\def\Curvesstack{\mathcal{C}\!\mathit{urves}}
\def\Polarizedstack{\mathcal{P}\!\mathit{olarized}}
\def\Complexesstack{\mathcal{C}\!\mathit{omplexes}}
% \Pic is the operator that assigns to X its picard group, usage \Pic(X)
% \Picardstack_{X/B} denotes the Picard stack of X over B
% \Picardfunctor_{X/B} denotes the Picard functor of X over B
\def\Pic{\mathop{\mathrm{Pic}}\nolimits}
\def\Picardstack{\mathcal{P}\!\mathit{ic}}
\def\Picardfunctor{\mathrm{Pic}}
\def\Deformationcategory{\mathcal{D}\!\mathit{ef}}
\usepackage{cleveref}
\crefname{originalreference}{original reference}{original references}
\Crefname{originalreference}{Original reference}{Original references}
\makeatletter
\newlabel{categories-lemma-inertia-fibred-category}{{categories-lemma-inertia-fibred-category}{source}{}{sourcepointer-1}{}}
\newlabel{categories-lemma-inertia-fibred-category@cref}{{[originalreference][0][]categories-lemma-inertia-fibred-category}{[1][0][]source}{}{}{}}
\makeatother
\crefname{originalreference}{original reference}{original references}
\Crefname{originalreference}{Original reference}{Original references}
\makeatletter
\newlabel{categories-lemma-relative-inertia-as-fibre-product}{{categories-lemma-relative-inertia-as-fibre-product}{source}{}{sourcepointer-2}{}}
\newlabel{categories-lemma-relative-inertia-as-fibre-product@cref}{{[originalreference][0][]categories-lemma-relative-inertia-as-fibre-product}{[1][0][]source}{}{}{}}
\makeatother
\crefname{originallemma}{original lemma}{original lemmas}
\Crefname{originallemma}{Original lemma}{Original lemmas}
\makeatletter
\newlabel{lemma-isom-locally-finite-type}{{lemma-isom-locally-finite-type}{source}{}{sourcepointer-3}{}}
\newlabel{lemma-isom-locally-finite-type@cref}{{[originallemma][0][]lemma-isom-locally-finite-type}{[1][0][]source}{}{}{}}
\makeatother
\crefname{originalreference}{original reference}{original references}
\Crefname{originalreference}{Original reference}{Original references}
\makeatletter
\newlabel{spaces-morphisms-lemma-permanence-finite-type}{{spaces-morphisms-lemma-permanence-finite-type}{source}{}{sourcepointer-4}{}}
\newlabel{spaces-morphisms-lemma-permanence-finite-type@cref}{{[originalreference][0][]spaces-morphisms-lemma-permanence-finite-type}{[1][0][]source}{}{}{}}
\makeatother
\title{The inertia stack and the relative inertia stack are algebraic and locally of finite type over the stack}
\author{The Stacks Project Authors}
\date{stacks-project:stacks-morphisms.tex:local-snapshot-20261012; GNU FDL 1.2 or later}
\begin{document}
\maketitle
\noindent\textit{Editorial note.} All delivered mathematical text is the author's own, copied from the pinned Stacks Project chapter source sources/stacks-algebraic-stacks-morphisms/pinned-original-source.tex; no statement or proof step was rewritten, reordered or removed.\par\medskip
\noindent\textit{Editorial note.} Preamble provenance: the Stacks Project keeps one single shared preamble for every chapter; this card takes it directly from the staged sources/stacks-algebraic-stacks-morphisms/preamble.tex as three verbatim spans (lines 8-12, 14-18 and 154-217), so no derived file was written and nothing was added or redefined. The wrapper class and the base packages are supplied by the assembler.\par\medskip
\noindent\textit{Editorial note.} Bibliography provenance: neither the counted statement nor its proof contains a citation (this lane rescans the whole span before writing the card), so the card carries no bibliography and no entry was reconstructed.\par\medskip
\noindent\textit{Editorial note.} Selection provenance: this card was opened under the file-based control contract (CONTROL.md), with the mechanical pre-card gate run on its statement and proof interval before the card was written.\par\medskip
\noindent\textit{Allowed standard prerequisites (editorial).}\begin{itemize}
\item The equivalence describing the inertia stack and the relative inertia stack as self fibre products of the diagonal. Invoked at the start of the counted proof (source lines 1239-1246) as a cross-chapter pointer. Shows that the inertia stacks are algebraic stacks.
\item The statement that the isomorphism space of two objects is an algebraic space locally of finite type. Invoked inside the counted proof (source lines 1256-1259). Makes the pullback of the inertia a space locally of finite type.
\item The description of the relative inertia as a fibre product of isomorphism spaces. Invoked inside the counted proof (source lines 1268-1271) as a cross-chapter pointer. Reduces the relative case to a comparison of two isomorphism spaces.
\item The permanence statement for locally of finite type morphisms of algebraic spaces. Invoked near the end of the counted proof (source lines 1274-1278) as a cross-chapter pointer. Gives local finite type of the comparison morphism.
\end{itemize}\medskip
\ifdefined\hypertarget\hypertarget{sourcepointer-1}{}\fi
\noindent\textit{Source reference.} Original reference categories-lemma-inertia-fibred-category: sources/stacks-algebraic-stacks-morphisms/pinned-original-source.tex, source lines 1240--1240. Navigation pointer only: the counted text refers to the Stacks item categories-lemma-inertia-fibred-category, which belongs to another chapter of the project and is not part of this card. The pointer carries the target's own label as its number because this chapter's staged snapshot does not contain the other chapter's tag table; its content is not reproduced here.\par\medskip
\ifdefined\hypertarget\hypertarget{sourcepointer-2}{}\fi
\noindent\textit{Source reference.} Original reference categories-lemma-relative-inertia-as-fibre-product: sources/stacks-algebraic-stacks-morphisms/pinned-original-source.tex, source lines 1282--1282. Navigation pointer only: the counted text refers to the Stacks item categories-lemma-relative-inertia-as-fibre-product, which belongs to another chapter of the project and is not part of this card. The pointer carries the target's own label as its number because this chapter's staged snapshot does not contain the other chapter's tag table; its content is not reproduced here.\par\medskip
\ifdefined\hypertarget\hypertarget{sourcepointer-3}{}\fi
\noindent\textit{Source reference.} Original lemma lemma-isom-locally-finite-type: sources/stacks-algebraic-stacks-morphisms/pinned-original-source.tex, source lines 278--278. Navigation pointer only: the counted text refers to the Stacks item lemma-isom-locally-finite-type, which is an existing labeled result of this same chapter and is not part of this card. The pointer carries the target's own label as its number and its definition line as its span; its content is not reproduced here.\par\medskip
\ifdefined\hypertarget\hypertarget{sourcepointer-4}{}\fi
\noindent\textit{Source reference.} Original reference spaces-morphisms-lemma-permanence-finite-type: sources/stacks-algebraic-stacks-morphisms/pinned-original-source.tex, source lines 307--307. Navigation pointer only: the counted text refers to the Stacks item spaces-morphisms-lemma-permanence-finite-type, which belongs to another chapter of the project and is not part of this card. The pointer carries the target's own label as its number because this chapter's staged snapshot does not contain the other chapter's tag table; its content is not reproduced here.\par\medskip
\begin{lemma}
\label{lemma-inertia}
Let $\mathcal{X}$ be an algebraic stack. Then the inertia stack
$\mathcal{I}_\mathcal{X}$ is an algebraic stack as well.
The morphism
$$
\mathcal{I}_\mathcal{X} \longrightarrow \mathcal{X}
$$
is representable by algebraic spaces and locally of finite type.
More generally, let $f : \mathcal{X} \to \mathcal{Y}$ be a morphism
of algebraic stacks. Then the relative inertia
$\mathcal{I}_{\mathcal{X}/\mathcal{Y}}$ is an algebraic stack and the
morphism
$$
\mathcal{I}_{\mathcal{X}/\mathcal{Y}} \longrightarrow \mathcal{X}
$$
is representable by algebraic spaces and locally of finite type.
\end{lemma}

\begin{proof}
By
Categories, Lemma \ref{categories-lemma-inertia-fibred-category}
there are equivalences
$$
\mathcal{I}_\mathcal{X} \to
\mathcal{X} \times_{\Delta, \mathcal{X} \times_S \mathcal{X}, \Delta}
\mathcal{X}
\quad\text{and}\quad
\mathcal{I}_{\mathcal{X}/\mathcal{Y}} \to
\mathcal{X}
\times_{\Delta, \mathcal{X} \times_\mathcal{Y} \mathcal{X}, \Delta}
\mathcal{X}
$$
which shows that the inertia stacks are algebraic stacks.
Let $T \to \mathcal{X}$ be a morphism given by
the object $x$ of the fibre category of $\mathcal{X}$ over $T$.
Then we get a $2$-fibre product square
$$
\xymatrix{
\mathit{Isom}_\mathcal{X}(x, x) \ar[d] \ar[r] &
\mathcal{I}_\mathcal{X} \ar[d] \\
T \ar[r]^x & \mathcal{X}
}
$$
This follows immediately from the definition of $\mathcal{I}_\mathcal{X}$.
Since $\mathit{Isom}_\mathcal{X}(x, x)$ is always an algebraic space
locally of finite type over $T$ (see
Lemma \ref{lemma-isom-locally-finite-type})
we conclude that $\mathcal{I}_\mathcal{X} \to \mathcal{X}$ is representable
by algebraic spaces and locally of finite type. Finally, for
the relative inertia we get
$$
\vcenter{
\xymatrix{
\mathit{Isom}_\mathcal{X}(x, x) \ar[d] &
K \ar[l] \ar[d] \ar[r] &
\mathcal{I}_{\mathcal{X}/\mathcal{Y}} \ar[d] \\
\mathit{Isom}_\mathcal{Y}(f(x), f(x)) &
T \ar[l]_-e \ar[r]^x & \mathcal{X}
}
}
$$
with both squares $2$-fibre products. This follows from
Categories, Lemma \ref{categories-lemma-relative-inertia-as-fibre-product}.
The left vertical arrow is a morphism of algebraic spaces locally of finite
type over $T$, and hence is locally of finite type, see
Morphisms of Spaces,
Lemma \ref{spaces-morphisms-lemma-permanence-finite-type}.
Thus $K$ is an algebraic space and $K \to T$ is locally of finite type.
This proves the assertion on the relative inertia.
\end{proof}

\end{document}
